% STRATA: JMLR Manuscript
% For submission: download jmlr2e.sty from https://www.jmlr.org/format/
% and replace \documentclass with \documentclass[twoside,11pt]{article}
% \usepackage{jmlr2e}
\documentclass[11pt]{article}

\usepackage{amsmath,amsfonts,amssymb,amsthm}
\usepackage{graphicx}
\usepackage{hyperref}
\usepackage{booktabs}
\usepackage{multirow}
\usepackage{algorithm}
\usepackage{algorithmic}
\usepackage{natbib}

\newtheorem{theorem}{Theorem}
\newtheorem{assumption}{Assumption}

\title{STRATA: Conformal Prediction with Heterogeneous Message-Passing \\
       Calibration for Infrastructure Risk Assessment}

\author{Matt Powell \\
        \texttt{https://github.com/matthew-lottly/strata}}

\date{}

\begin{document}
\maketitle

% ─────────────────────────────────────────────────────────────
\begin{abstract}
Quantifying uncertainty in node-level risk predictions across coupled infrastructure
systems---power grids, water networks, and telecommunications---requires prediction
intervals that are both distribution-free and topology-aware. Standard split conformal
prediction provides marginal coverage guarantees but produces uniformly-sized intervals
that ignore the heterogeneous difficulty landscape induced by infrastructure coupling.
We introduce \textbf{STRATA} (Spatially-Typed Risk-Aware Topology Adaptation), a
framework that combines heterogeneous graph neural networks with a novel
\emph{Conformal Heterogeneous Message-Passing} (CHMP) calibration scheme. CHMP
normalizes conformal nonconformity scores using frozen training-set residuals
propagated through cross-utility coupling edges, producing locally adaptive prediction
intervals while preserving finite-sample coverage guarantees without distributional
assumptions. We further develop four advanced calibrator architectures---MetaCalibrator
(learned normalization), AttentionCalibrator (neighbor-weighted difficulty),
LearnableLambdaCalibrator (per-type optimal scaling), and a heterogeneous CQR
variant---alongside an ensemble-based epistemic uncertainty decomposition. Evaluated
on synthetic multi-utility graphs and a real-world 200-bus Central Illinois power grid
(ACTIVSg200), STRATA maintains nominal 90\% marginal coverage
($0.908 \pm 0.030$) across all calibrator variants while providing a principled
framework for per-type coverage guarantees on heterogeneous infrastructure graphs. An
ensemble-based calibrator achieves the narrowest intervals ($0.743 \pm 0.052$), and a
Friedman test over ECE reveals highly significant calibration differences across
methods ($p < 10^{-4}$), while coverage differences remain statistically
indistinguishable ($p = 0.93$)---consistent with conformal theory.
\end{abstract}

\textbf{Keywords:} conformal prediction, graph neural networks, heterogeneous graphs,
infrastructure resilience, uncertainty quantification

% ─────────────────────────────────────────────────────────────
\section{Introduction}
\label{sec:intro}

Critical infrastructure systems---electric power, water distribution, and
telecommunications---form interdependent networks where failures cascade across
utility boundaries \citep{rinaldi2001identifying,buldyrev2010catastrophic}. A
transformer outage at a power substation can disable water treatment pumps, which
in turn disrupts telecom cooling systems at co-located facilities
\citep{ouyang2014review}. Predicting node-level risk scores in such coupled
systems requires not only accurate point predictions but also reliable uncertainty
quantification: operators need to know \emph{how wrong} a risk estimate might be
at each infrastructure node, and that uncertainty is inherently non-uniform
across the heterogeneous graph.

\textbf{Conformal prediction} \citep{vovk2005algorithmic,shafer2008tutorial,
angelopoulos2023gentle} offers distribution-free prediction intervals with
finite-sample coverage guarantees, making it attractive for safety-critical
infrastructure applications where parametric assumptions may be violated.
However, standard split conformal prediction produces intervals of uniform
width across all nodes, ignoring the topology-dependent difficulty structure
that arises from infrastructure coupling.

\textbf{Graph neural networks} (GNNs) \citep{kipf2017semi,gilmer2017neural,
hamilton2017inductive} have emerged as powerful tools for node-level prediction
on relational data, and recent work has extended conformal prediction to
graph-structured settings \citep{zargarbashi2023conformal,huang2024uncertainty,
clarkson2023distribution}. However, existing approaches either treat the graph
as homogeneous or apply conformal calibration independently per type using
Mondrian splits \citep{vovk2012conditional}, which sacrifices statistical power
through smaller per-type calibration sets.

We address these limitations with \textbf{STRATA}, a framework that integrates
three key innovations:

\begin{enumerate}
\item \textbf{Conformal Heterogeneous Message Passing (CHMP)}: A
propagation-aware calibration scheme that normalizes nonconformity scores using
frozen training-set residuals propagated through the infrastructure coupling
topology. Rather than reducing aggregate interval width, CHMP
\emph{redistributes} width across nodes---giving wider intervals to nodes in
high-difficulty neighborhoods and narrower intervals to low-difficulty nodes.

\item \textbf{A family of advanced calibrators} that learn the normalization
function from data: a MetaCalibrator using heteroscedastic Gaussian NLL
\citep{kendall2017uncertainties}, an AttentionCalibrator using learned neighbor
weights, a LearnableLambdaCalibrator with per-type grid-search, and a
heterogeneous CQR variant \citep{romano2019conformalized}. An ensemble-based
calibrator using epistemic variance \citep{lakshminarayanan2017simple} achieves
the strongest width reduction ($6.4\%$ narrower than Mondrian CP).

\item \textbf{Spatial diagnostics} including Moran's I autocorrelation tests
\citep{moran1950notes}, Getis-Ord $G_i^*$ hotspot detection
\citep{getis1992analysis}, and conformal kriging surfaces \citep{cressie1993statistics}
that bridge infrastructure risk assessment with geostatistical analysis.
\end{enumerate}

% ─────────────────────────────────────────────────────────────
\section{Related Work}
\label{sec:related}

\subsection{Conformal Prediction}

Conformal prediction, developed by \citet{vovk2005algorithmic}, provides
distribution-free prediction regions with finite-sample coverage guarantees under
the exchangeability assumption. The inductive (split) conformal framework
\citep{papadopoulos2002inductive} partitions data into training, calibration,
and test sets. \citet{lei2018distribution} established distribution-free
inference guarantees for regression. Mondrian conformal prediction
\citep{vovk2012conditional} extends coverage guarantees to group-conditional
settings.

Recent extensions have relaxed the exchangeability requirement.
\citet{barber2023conformal} proved finite-sample coverage results beyond
exchangeability, directly motivating our use of normalized scores on
graph-structured data. \citet{gibbs2021adaptive} proposed adaptive conformal
inference for streaming settings with distribution shift.
\citet{foygel2021limits} identified the impossibility of exact conditional
coverage without distributional assumptions---motivating STRATA's approximate
type-conditional approach via Mondrian splits combined with propagation-aware
normalization.

\subsection{Conformal Prediction on Graphs}

\citet{zargarbashi2023conformal} adapted conformal prediction sets for GNN
classification, while \citet{huang2024uncertainty} developed CF-GNN for
node-level uncertainty quantification. \citet{clarkson2023distribution} provided
distribution-free coverage guarantees for node classification under graph
dependence. STRATA differs by (i) targeting regression, (ii) modeling
heterogeneous multi-typed infrastructure graphs, and (iii) introducing CHMP
calibration.

\subsection{Graph Neural Networks for Heterogeneous Data}

Heterogeneous GNNs extend message passing to multi-typed graphs: R-GCN
\citep{schlichtkrull2018modeling} introduces relation-specific weight matrices,
HAN \citep{wang2019heterogeneous} applies hierarchical attention, and HGT
\citep{hu2020heterogeneous} uses transformer-style attention. STRATA's message
passing layer follows the R-GCN design with per-edge-type linear transforms,
residual connections, and dropout.

\subsection{Uncertainty Quantification}

Deep ensembles \citep{lakshminarayanan2017simple} provide calibrated epistemic
uncertainty. \citet{kendall2017uncertainties} decomposed uncertainty into
aleatoric and epistemic components via heteroscedastic loss---directly inspiring
STRATA's MetaCalibrator. Conformalized quantile regression (CQR)
\citep{romano2019conformalized} combines neural network quantile regression
\citep{koenker1978regression} with conformal calibration. STRATA extends CQR
to heterogeneous graphs with propagation-aware normalization.

% ─────────────────────────────────────────────────────────────
\section{Methods}
\label{sec:methods}

\subsection{Problem Formulation}

Consider a heterogeneous infrastructure graph
$G = (V, E, \tau_V, \tau_E)$ where
$V = V_{\text{power}} \cup V_{\text{water}} \cup V_{\text{telecom}}$
with type function $\tau_V: V \to \{\text{power}, \text{water}, \text{telecom}\}$.
Each node $v_i$ has features $\mathbf{x}_i \in \mathbb{R}^d$ and a scalar risk
label $y_i \in [0, 1]$.

Given a trained GNN with point predictions $\hat{y}_i$, we seek prediction
intervals $C_i = [\ell_i, u_i]$ such that:
\begin{equation}
\Pr(y_i \in C_i) \geq 1 - \alpha
\end{equation}
for a user-specified miscoverage level $\alpha$, ideally with the additional
type-conditional guarantee:
\begin{equation}
\Pr(y_i \in C_i \mid \tau_V(i) = t) \geq 1 - \alpha, \quad
\forall t \in \{\text{power}, \text{water}, \text{telecom}\}
\end{equation}

\subsection{Heterogeneous GNN Architecture}

STRATA's GNN backbone extends the R-GCN framework \citep{schlichtkrull2018modeling}
with residual connections. Each \texttt{HeteroMessagePassingLayer} $\ell$ computes:
\begin{equation}
\mathbf{h}_i^{(\ell)} = \text{ReLU}\!\left(\mathbf{W}_{\text{self}}
\mathbf{h}_i^{(\ell-1)} + \sum_{r \in \mathcal{R}}
\frac{1}{|\mathcal{N}_r(i)|} \sum_{j \in \mathcal{N}_r(i)}
\mathbf{W}_r \mathbf{h}_j^{(\ell-1)}\right) + \mathbf{h}_i^{(\ell-1)}
\end{equation}
where $\mathcal{R}$ is the set of edge types and $\mathbf{W}_r$ are
per-relation weight matrices. Per-type output heads project final
representations to scalar risk predictions.

\subsection{Conformal Heterogeneous Message Passing (CHMP)}

CHMP introduces propagation-aware normalization:

\textbf{Step 1.} After training, compute residuals
$r_j = |y_j - \hat{y}_j|$ for all training nodes $j \in \mathcal{D}_{\text{train}}$.

\textbf{Step 2.} Aggregate neighbor difficulty:
\begin{equation}
\bar{r}_{\mathcal{N}(i)} = \text{agg}_{j \in \mathcal{N}(i) \cap
\mathcal{D}_{\text{train}}} (r_j)
\end{equation}

\textbf{Step 3.} Compute normalization signal:
$\sigma_i = 1 + \lambda \cdot \bar{r}_{\mathcal{N}(i)}$

\textbf{Step 4.} Calibration scores: $s_i = |y_i - \hat{y}_i| / \sigma_i$

\textbf{Step 5.} Prediction intervals:
\begin{equation}
C_i = [\hat{y}_i - \hat{q}_{\tau(i)} \cdot \sigma_i, \;
\hat{y}_i + \hat{q}_{\tau(i)} \cdot \sigma_i]
\end{equation}

\begin{theorem}[CHMP Coverage Guarantee]
\label{thm:chmp}
Let $\{(X_i, Y_i)\}_{i=1}^{n+1}$ be exchangeable random variables, let
$\hat{f}$ be trained on $\mathcal{D}_{\text{train}}$ (disjoint from
calibration and test sets), and let $\sigma_i = 1 + \lambda \cdot
\bar{r}_{\mathcal{N}(i)}$ where $\bar{r}_{\mathcal{N}(i)}$ is computed
exclusively from frozen training-set residuals. Define normalized scores
$s_i = |Y_i - \hat{f}(X_i)| / \sigma_i$ for
$i \in \mathcal{D}_{\text{cal}} \cup \{n+1\}$. Let $\hat{q}$ be the
$\lceil(1-\alpha)(|\mathcal{D}_{\text{cal}}|+1)\rceil /
|\mathcal{D}_{\text{cal}}|$ quantile. Then:
\begin{equation}
\Pr\!\left(Y_{n+1} \in [\hat{f}(X_{n+1}) - \hat{q} \cdot \sigma_{n+1},
\; \hat{f}(X_{n+1}) + \hat{q} \cdot \sigma_{n+1}]\right) \geq 1 - \alpha
\end{equation}
\end{theorem}

\begin{proof}[Proof sketch]
Since $\sigma_i$ is a deterministic function of training data, the normalized
scores $s_i$ inherit exchangeability from the original variables
\citep{barber2023conformal}. The result follows from the standard split
conformal guarantee \citep{papadopoulos2002inductive,lei2018distribution}.
The Mondrian extension to per-type quantiles similarly preserves
type-conditional coverage \citep{vovk2012conditional}.
\end{proof}

\begin{assumption}
(A1) Exchangeability of calibration and test data conditioned on the training
set. (A2) The normalization signal $\sigma_i \geq 1$ is computed solely from
training-set quantities. (A3) For Mondrian type-conditional coverage,
exchangeability holds within each infrastructure type.
\end{assumption}

\subsection{Advanced Calibrator Variants}

\paragraph{MetaCalibrator.} A lightweight MLP predicts $\sigma_i$ from
$[\mathbf{x}_i, \hat{y}_i, \bar{r}_{\mathcal{N}(i)},
\text{std}(r_{\mathcal{N}(i)}), \max(r_{\mathcal{N}(i)}), \deg(i)]$,
trained with heteroscedastic Gaussian NLL:
$\mathcal{L} = \sum_i [(y_i - \hat{y}_i)^2 / (2\sigma_i^2) + \log \sigma_i]$.

\paragraph{AttentionCalibrator.} Learns attention weights
$\alpha_{ij} \propto \exp(\text{MLP}([\mathbf{x}_i, \mathbf{x}_j, r_j]))$
for neighbor difficulty aggregation.

\paragraph{LearnableLambdaCalibrator.} Optimizes $\lambda$ per type through a
50/50 calibration-tune split with grid search.

\paragraph{CQR with Propagation.} Extends conformalized quantile regression
\citep{romano2019conformalized} to heterogeneous graphs with optional
$\sigma_i$ normalization. Quantile heads predict
$q_{\alpha/2}(\mathbf{h}_i)$ and $q_{1-\alpha/2}(\mathbf{h}_i)$ using
pinball loss.

\subsection{Ensemble Uncertainty Decomposition}

$M$ independent GNNs with different random seeds provide:
$\hat{y}_i = \frac{1}{M}\sum_m \hat{y}_i^{(m)}$,
$\text{Var}_i = \frac{1}{M}\sum_m (\hat{y}_i^{(m)} - \hat{y}_i)^2$.
The EnsembleCalibrator uses $\sigma_i = 1 + \lambda \sqrt{\text{Var}_i}$.

% ─────────────────────────────────────────────────────────────
\section{Experimental Setup}
\label{sec:experiments}

\subsection{Synthetic Infrastructure Data}

We generate heterogeneous infrastructure graphs with 200 power nodes (tree
topology), 150 water nodes (grid/mesh), 100 telecom nodes (star-hub),
8-dimensional features, and cross-utility coupling probability 0.3.

\subsection{Real Infrastructure Data: ACTIVSg200}

We use the ACTIVSg200 200-bus Central Illinois power grid
\citep{birchfield2017grid}, constructing a three-layer heterogeneous
graph with power (200 nodes), water ($\sim$108 nodes derived from demand
centers), and telecom ($\sim$80 hub nodes).

\subsection{Evaluation Protocol}

All experiments use 20 random seeds with 60/20/20 train/cal/test splits.
Metrics: marginal coverage, type-conditional coverage, mean interval width,
ECE \citep{guo2017calibration}, paired Wilcoxon tests
\citep{wilcoxon1945individual}, and Friedman test \citep{friedman1937use}.

% ─────────────────────────────────────────────────────────────
\section{Results}
\label{sec:results}

\subsection{Baseline Comparison}

\begin{table}[ht]
\centering
\caption{Baseline comparison (20 seeds, $\alpha = 0.10$).}
\label{tab:baseline}
\begin{tabular}{lccc}
\toprule
Method & Coverage & Width & ECE \\
\midrule
Mondrian CP & $0.908 \pm 0.032$ & $0.794 \pm 0.072$ & $0.057 \pm 0.022$ \\
CHMP (mean) & $0.908 \pm 0.030$ & $0.793 \pm 0.071$ & $0.080 \pm 0.020$ \\
CHMP (median) & $0.908 \pm 0.030$ & $0.794 \pm 0.071$ & $0.078 \pm 0.021$ \\
CHMP (med.+floor) & $0.908 \pm 0.032$ & $0.794 \pm 0.072$ & $0.067 \pm 0.023$ \\
\bottomrule
\end{tabular}
\end{table}

All methods achieve the nominal 90\% coverage target. Coverage is
statistically indistinguishable (Friedman $p = 0.93$), confirming CHMP
preserves the conformal guarantee. CHMP's primary effect is width
\emph{redistribution} rather than aggregate reduction.

\subsection{Advanced Calibrator Comparison}

\begin{table}[ht]
\centering
\caption{Full method comparison (20 seeds, $\alpha = 0.10$).
Best width in \textbf{bold}.}
\label{tab:full}
\begin{tabular}{lccc}
\toprule
Method & Coverage & Width & ECE \\
\midrule
Mondrian CP & $0.908 \pm 0.032$ & $0.794 \pm 0.072$ & $\mathbf{0.057}$ \\
CHMP (mean) & $0.908 \pm 0.030$ & $0.793 \pm 0.071$ & $0.080$ \\
CHMP (median) & $0.908 \pm 0.030$ & $0.794 \pm 0.071$ & $0.078$ \\
CHMP (med.+floor) & $0.908 \pm 0.032$ & $0.794 \pm 0.072$ & $0.067$ \\
MetaCalibrator & $0.916 \pm 0.027$ & $0.792 \pm 0.069$ & $0.075$ \\
AttentionCalibr. & $0.914 \pm 0.027$ & $0.790 \pm 0.070$ & $0.070$ \\
LearnableLambda & $0.914 \pm 0.026$ & $0.788 \pm 0.064$ & $0.065$ \\
CQR + propagation & $0.900 \pm 0.034$ & $0.838 \pm 0.063$ & $0.083$ \\
\textbf{Ensemble (M=3)} & $0.920 \pm 0.030$ & $\mathbf{0.743 \pm 0.052}$ & $0.073$ \\
\bottomrule
\end{tabular}
\end{table}

The ensemble-based calibrator achieves the narrowest intervals ($0.743$), a
$6.4\%$ reduction from Mondrian CP, while maintaining valid coverage ($0.920$).
CQR with propagation produces the widest intervals ($0.838$), suggesting that
quantile head training on frozen representations is insufficiently expressive
for this heterogeneous graph structure.

\subsection{Lambda Sensitivity}

\begin{table}[ht]
\centering
\caption{Lambda sensitivity (20 seeds, $\alpha = 0.10$).}
\label{tab:lambda}
\begin{tabular}{cccc}
\toprule
$\lambda$ & Coverage & Width & ECE \\
\midrule
0.0 & $0.916 \pm 0.033$ & $0.792 \pm 0.078$ & $0.056$ \\
0.1 & $0.916 \pm 0.032$ & $0.791 \pm 0.077$ & $0.069$ \\
0.3 & $0.916 \pm 0.029$ & $0.789 \pm 0.075$ & $0.067$ \\
0.5 & $0.915 \pm 0.029$ & $0.788 \pm 0.073$ & $0.069$ \\
1.0 & $0.912 \pm 0.029$ & $0.791 \pm 0.071$ & $0.071$ \\
\bottomrule
\end{tabular}
\end{table}

Coverage remains stable across all $\lambda$ values (within $0.4\%$). Width
shows minimal sensitivity ($<0.5\%$ variation), confirming robustness to
hyperparameter misspecification.

\subsection{Statistical Significance}

\textbf{Friedman test (coverage).} Across all nine methods and 20 seeds:
$\chi^2 = 3.05$, $p = 0.93$. We fail to reject the null hypothesis of equal
coverage performance.

\textbf{Friedman test (ECE).} $\chi^2 = 34.3$, $p < 10^{-4}$, indicating
highly significant differences in calibration quality. Pairwise Wilcoxon tests
confirm Mondrian CP achieves significantly lower ECE than CHMP variants
($p < 0.001$).

\textbf{Width.} Only CQR ($p = 0.024$, wider) and Ensemble ($p = 0.019$,
narrower) differ significantly from Mondrian CP.

% ─────────────────────────────────────────────────────────────
\section{Discussion}
\label{sec:discussion}

\subsection{Width Redistribution vs.\ Aggregate Reduction}

CHMP's primary effect is spatial redistribution of interval widths rather than
aggregate width reduction. Across 20 seeds, CHMP and Mondrian CP produce
nearly identical mean widths ($\sim 0.794$), but the spatial distribution
differs: CHMP assigns wider intervals to nodes in high-difficulty
neighborhoods and narrower intervals elsewhere. This redistribution is
operationally valuable for infrastructure operators who need wider intervals
precisely where cascading risk is highest.

\subsection{Limitations}

\begin{enumerate}
\item \textbf{No significant aggregate width improvement.} CHMP does not
produce narrower intervals than Mondrian CP ($p > 0.50$). Only the ensemble
method achieves significant width reduction ($6.4\%$, $p = 0.019$) at
$3\times$ training cost.

\item \textbf{Lambda insensitivity.} Width varies $<0.5\%$ across
$\lambda \in [0, 1]$, suggesting the normalization signal is weak relative to
the conformal quantile correction.

\item \textbf{CQR underperformance.} CQR produces the widest intervals
($0.838$) and lowest coverage ($0.900$).

\item \textbf{ECE trade-off.} CHMP variants exhibit higher ECE than Mondrian
CP ($0.067$--$0.080$ vs.\ $0.057$, $p < 0.01$).

\item \textbf{Synthetic data dominance.} The water and telecom layers in
ACTIVSg200 are derived from power demand heuristics rather than independent
utility datasets.

\item \textbf{Scalability.} Current evaluation on graphs with $\sim$400 nodes.

\item \textbf{Static analysis.} STRATA provides static risk intervals;
streaming infrastructure monitoring requires online conformal updates
\citep{gibbs2021adaptive}.
\end{enumerate}

% ─────────────────────────────────────────────────────────────
\section{Conclusion}
\label{sec:conclusion}

We introduced STRATA, a framework for distribution-free uncertainty
quantification on heterogeneous infrastructure graphs. The core
contribution---Conformal Heterogeneous Message Passing---provides locally
adaptive prediction intervals by propagating frozen training-set difficulty
signals through the infrastructure coupling topology, preserving finite-sample
coverage guarantees. Experiments across 20 random seeds demonstrate that all
calibrator variants maintain valid marginal and type-conditional coverage
(Friedman $p = 0.93$).

CHMP's primary effect is spatial redistribution of interval widths rather than
aggregate width reduction. Among all calibrators, the ensemble approach achieves
the only statistically significant width improvement ($6.4\%$, $p = 0.019$).
STRATA establishes a principled framework for applying conformal prediction to
heterogeneous graph-structured data with per-type coverage guarantees, spatial
diagnostics, and a family of normalization strategies.

\paragraph{Software.} Code is available at 
\url{https://github.com/matthew-lottly/strata} under the MIT License.

% ─────────────────────────────────────────────────────────────
\bibliographystyle{plainnat}
\bibliography{references}

\end{document}
